\documentclass{article}
\usepackage{amsmath,amssymb}
\usepackage{xurl}
\usepackage[hidelinks]{hyperref}
% Shared commands for notes in docs/paper-gaps/.

\DeclareUnicodeCharacter{2080}{\ifmmode{}_{0}\else\textsubscript{0}\fi}
\DeclareUnicodeCharacter{2081}{\ifmmode{}_{1}\else\textsubscript{1}\fi}
\DeclareUnicodeCharacter{2082}{\ifmmode{}_{2}\else\textsubscript{2}\fi}
\DeclareUnicodeCharacter{2099}{\ifmmode{}_{n}\else\textsubscript{n}\fi}

\newcommand{\Lean}{\textsc{Lean}}
\ExplSyntaxOn
\NewDocumentCommand{\leanid}{m}{
  \ifmmode
    \textsf{\detokenize{#1}}
  \else
    \group_begin:
    \urlstyle{sf}
    \tl_set:Nn \l_tmpa_tl {#1}
    \tl_replace_all:Nnn \l_tmpa_tl {\_} {_}
    \exp_args:NV \path \l_tmpa_tl
    \group_end:
  \fi
}
\ExplSyntaxOff
\providecommand{\tr}{\operatorname{tr}}
\newcommand{\tnleanIssueBase}{https://github.com/LionSR/TNLean/issues/}
\newcommand{\tnleanPrBase}{https://github.com/LionSR/TNLean/pull/}
\newcommand{\tnleanDocBase}{https://sirui-lu.com/TNLean/docs/}
\newcommand{\ghissue}[1]{\href{\tnleanIssueBase#1}{GitHub issue \##1}}
\newcommand{\ghpr}[1]{\href{\tnleanPrBase#1}{PR \##1}}
\newcommand{\doclink}[1]{\href{\tnleanDocBase#1.html}{\path{#1}}}

% Dirac notation and the identity operator, as in blueprint/src/macros/common.tex.
% \providecommand keeps note-local definitions authoritative.
\usepackage{dsfont}
\providecommand{\ket}[1]{|#1\rangle}
\providecommand{\bra}[1]{\langle #1|}
\providecommand{\braket}[2]{\langle #1 | #2 \rangle}
\providecommand{\ketbra}[2]{|#1\rangle\!\langle #2|}
\providecommand{\Id}{\mathds{1}}
\providecommand{\one}{\mathds{1}}
\providecommand{\C}{\mathbb{C}}
\providecommand{\MN}[1]{M_{#1}(\C)}
\providecommand{\MD}{M_D(\C)}

% External referencing for paper sources.  The build helper writes sanitized
% auxiliary files under build/paper-aux/ and excludes duplicated source labels.
\usepackage{xr}

% Import a generated paper auxiliary file with a caller-supplied label prefix.
% The candidate paths support builds from docs/paper-gaps/, the repository root,
% and build/paper-aux/ itself.
\newcommand{\importpaperaux}[2]{%
  \IfFileExists{../../build/paper-aux/#2.aux}{%
    \externaldocument[#1]{../../build/paper-aux/#2}%
  }{%
    \IfFileExists{build/paper-aux/#2.aux}{%
      \externaldocument[#1]{build/paper-aux/#2}%
    }{%
      \IfFileExists{#2.aux}{%
        \externaldocument[#1]{#2}%
      }{}%
    }%
  }%
}

% General external reference macros
\newcommand{\paperref}[2]{\ref{#1-#2}}
\newcommand{\papereqref}[2]{(\ref{#1-#2})}

% CPSV16 source (1606.00608 / MPDO-22-12-17-2.tex) external document and shortcuts
\importpaperaux{cpsv16-}{1606.00608/MPDO-22-12-17-2-tnlean}
\newcommand{\cpsvref}[1]{\paperref{cpsv16}{#1}}
\newcommand{\cpsveqref}[1]{\papereqref{cpsv16}{#1}}

% Verdict marker: \gapnote{<kind>}{<status>}. Kind is the policy
% classification (clarification, local-correction, scope-restriction,
% unfaithful, false-source, open-gap); status is open, wip (elimination
% underway), resolved, or historical. Machine-read by the site index;
% prints a verdict line.
\newcommand{\gapnote}[2]{\par\noindent\textbf{Verdict:} #1 (#2).\par}

\title{Kitaev Chain: the Bosonic Contraction of the Fixed-Point Matrices}
\author{}
\date{2026-09-25}
\begin{document}
\maketitle
\gapnote{scope-restriction}{open}

\section{What the sources print}
Appendix~A of \cite{Cirac2021Matrix}, under ``The Kitaev chain'', describes the
non-trivial fixed point of the Kitaev chain \cite{Kitaev2001Majorana} as a
graded (fermionic) tensor network with matrices
\begin{align}
    A^0=\begin{pmatrix}1&0\\0&1\end{pmatrix},\qquad
    A^1=\begin{pmatrix}0&1\\-1&0\end{pmatrix},
    \label{eq:kitaev_tensor}
\end{align}
and a twist $Y=A^1$ at the boundary.\footnote{Source traceability:
    \path{Papers/2011.12127/TN-Review-main.tex}, lines 2526--2532.}
The state it refers to is
\begin{align}
    \ket\psi=\sum_{i_1,\ldots,i_N}\tr\bigl(YA^{i_1}\cdots A^{i_N}\bigr)
    \ket{i_1}\otimes_g\ket{i_2}\otimes_g\cdots\otimes_g\ket{i_N},
    \label{eq:kitaev_state}
\end{align}
where $\otimes_g$ is the graded tensor product; the review states that this
state has odd parity and that, contrary to the bosonic case, the description is
injective.\footnote{Source traceability:
    \path{Papers/2011.12127/TN-Review-main.tex}, lines 474--478.}
The matrices and the extra $y$ in the periodic trace are those of
\cite{Bultinck2017FermionicMPS}, which also notes that $Y$ commutes with every
$A^i$, so that the matrices read as a bosonic matrix product state are
reducible.\footnote{Source traceability:
    \path{References/1610.07849/source/FermionicMPS.tex}, lines 264--289 and 372.}

\section{The restriction}
The graded tensor product, the sign rules of the fermionic virtual
contraction, and the graded notion of injectivity are not formalized. Every
formal statement about this example concerns the ordinary contraction of the
matrices~\eqref{eq:kitaev_tensor}. This is enough for the coefficients
of~\eqref{eq:kitaev_state}: they are the ordinary traces $\tr(YA^{i_1}\cdots
A^{i_N})$, taken relative to the ordered graded product basis. The statements
proved are the following, with $k$ the number of letters equal to $1$ in the
word $i_1\cdots i_N$.\footnote{Results: \leanid{MPSTensor.kitaevTensor},
    \leanid{MPSTensor.kitaevTensor_coeff},
    \leanid{MPSTensor.kitaevTwistedCoeff_eq},
    \leanid{MPSTensor.kitaevTensor_not_isNormal}.}
\begin{enumerate}
    \item $(A^1)^2=-\mathbb 1$, $A^1=Y$, and $Y$ commutes with $A^0$ and $A^1$.
    \item $A^{i_1}\cdots A^{i_N}=Y^k$, hence
    $\tr(A^{i_1}\cdots A^{i_N})=2(-1)^{k/2}$ for even $k$ and $0$ for odd $k$,
    while $\tr(YA^{i_1}\cdots A^{i_N})=0$ for even $k$ and $2(-1)^{(k+1)/2}$
    for odd $k$. Neither depends on the length $N$ beyond $k$.
    \item The twisted amplitudes are supported exactly on odd $k$, the
    bosonic counterpart of the odd parity in~\eqref{eq:kitaev_state}; the
    untwisted amplitudes are supported exactly on even $k$.
    \item For every word, the $(0,0)$ and $(1,1)$ entries of the product agree
    and the $(0,1)$ and $(1,0)$ entries differ by a sign; these are the
    bosonic counterparts of the open-boundary closures discussed in
    \cite{Bultinck2017FermionicMPS}.
    \item The vector $(1,i)$ is a common eigenvector of $A^0$ and $A^1$, and
    the tensor is neither injective nor normal as a bosonic matrix product
    state. This is the ``bosonic case'' that the review contrasts with the
    fermionic one; the fermionic injectivity claim is not formalized.
\end{enumerate}

\section{Elimination plan}
The restriction is lifted once graded tensor networks are formalized: a
$\mathbb Z_2$-graded virtual space with parity operator $\operatorname{diag}(1,-1)$
under which $A^0$ is even and $A^1$ is odd, the graded contraction, and the
graded injectivity notion of \cite{Bultinck2017FermionicMPS}. The review's
claim that the graded description is injective would then be stated against
that notion.

\bibliographystyle{alpha}
\bibliography{references}
\end{document}
